\section*{الفصل \ref{ch:g10:chemistry-of-life} --- التركيب الكيميائي للكائنات الحية}

\begin{solution}{exo:g10:chemistry-of-life:1}
الأكسجين (نحو 65\%)، والكربون (18.5\%)، والهيدروجين (9.5\%)،
والنيتروجين (3.2\%): أي 96\% من كتلة الجسم مجتمعة.
\end{solution}

\begin{solution}{exo:g10:chemistry-of-life:2}
معدنية: الماء، وكلور الصوديوم، وفوسفات الكالسيوم، وثاني أكسيد الكربون
(مركب كربوني، لكنه يصنع خارج الحياة وليس له هيكل كربون--هيدروجين).
عضوية: الغلوكوز، وثلاثي الغليسيريد، \omterm{def:g10:chemistry-of-life:families}{والحمض الأميني}.
\end{solution}

\begin{solution}{exo:g10:chemistry-of-life:3}
الماء المفقود $12.0 - 0.6 = \qty{11.4}{g}$؛ و$11.4/12.0 = 0.95$: أي
أن الورقة 95\% ماء.
\end{solution}

\begin{solution}{exo:g10:chemistry-of-life:4}
(أ) محلول اليود على البذرة المسحوقة: أزرق أسود. (ب) كاشف البيوريه:
بنفسجي. (ج) اسحق البذرة على ورق: بقعة شفافة لا تجف؛ والسودان III
يلونها بالأحمر.
\end{solution}

\begin{solution}{exo:g10:chemistry-of-life:5}
\omterm{def:g10:chemistry-of-life:families}{الكربوهيدرات}: سكريات بسيطة مثل الغلوكوز. \omterm{def:g10:chemistry-of-life:families}{الدهون}: غليسرول وحموض دهنية.
\omterm{def:g10:chemistry-of-life:families}{البروتينات}: حموض أمينية (عشرون نوعًا). \omterm{def:g10:chemistry-of-life:families}{الحموض النووية}: نوكليوتيدات
(أربعة أنواع).
\end{solution}

\begin{solution}{exo:g10:chemistry-of-life:6}
الكربون: $18.5/0.02 \approx 900$ مرة أكثر تركزًا. والنيتروجين:
$3.2/0.002 = 1600$ مرة. فالمادة الحية ليست عينة من محيطها: إنها تجمع
عناصر نادرة قليلة بعوامل تبلغ الألف، وذلك توقيع كيميائي للحياة.
\end{solution}

\begin{solution}{exo:g10:chemistry-of-life:7}
الماء: $0.62 \times 70 = \qty{43.4}{kg}$. المادة الجافة: $70 - 43.4 =
\qty{26.6}{kg}$. \omterm{def:g10:chemistry-of-life:families}{البروتين}: $0.17 \times 70 = \qty{11.9}{kg}$.
\end{solution}

\begin{solution}{exo:g10:chemistry-of-life:8}
\omterm{def:g10:chemistry-of-life:families}{الكربوهيدرات} $60 \times 17 = \qty{1020}{kJ}$؛ \omterm{def:g10:chemistry-of-life:families}{والدهون} $12 \times 38 =
\qty{456}{kJ}$؛ \omterm{def:g10:chemistry-of-life:families}{والبروتينات} $8 \times 17 = \qty{136}{kJ}$؛ والمجموع
\qty{1612}{kJ}، أي نحو \qty{1600}{kJ}. وتقدم \omterm{def:g10:chemistry-of-life:families}{الدهون} $456/1612 \approx
28\%$ من الطاقة بنسبة 12\% فقط من الكتلة.
\end{solution}

\begin{solution}{exo:g10:chemistry-of-life:9}
الخبز: أزرق أسود --- فالدقيق نشاء في معظمه. الجبن: لا تغير في اللون
وراء صبغة اليود نفسها --- بروتينات ودهون، ولا نشاء. السكر: لا تغير
--- فالسكروز سكر بسيط لا نشاء؛ واليود لا يكشف إلا سلسلة النشاء
الطويلة.
\end{solution}

\begin{solution}{exo:g10:chemistry-of-life:10}
المادة الجافة عضوية في معظمها: جزيئات كربون--هيدروجين تتحد مع
الأكسجين فتحرر طاقة وثاني أكسيد كربون وماء. أما الرماد فهو \omterm{def:g10:chemistry-of-life:mineral-organic}{الأملاح
المعدنية}، وقد اتحدت بالأكسجين اتحادًا تامًا من قبل؛ فليس فيها ما
يحترق.
\end{solution}

\begin{solution}{exo:g10:chemistry-of-life:11}
يتفاعل كاشف فهلنغ مع زمرة كيميائية حرة تحملها جزيئات الغلوكوز
المفردة؛ أما في النشاء فوحدات الغلوكوز مرتبط بعضها ببعض وتلك الزمر
محبوسة داخل السلسلة، فلا يجد الكاشف شيئًا. وعلى العكس، يستقر اليود في
سلسلة النشاء الطويلة الملتفة، وهي سلسلة لا تكونها جزيئات الغلوكوز
المفردة. فكل كاشف يكشف بنية، لا عنصرًا.
\end{solution}

\begin{solution}{exo:g10:chemistry-of-life:12}
امتصت الحبة الماء امتصاصًا هائلًا. فعند 13\% تكون المدخرات صلبة
والإنزيمات عاجزة عن العمل: كيمياء الحبة مطفأة وهي تبقى ساكنة. والماء
يحل المدخرات، ويتيح للإنزيمات أن تعمل، وينقل النواتج إلى الرشيم
النامي: فالكيمياء لا تعود إلى الاشتغال إلا حين يوفر الماء.
\end{solution}

\begin{solution}{exo:g10:chemistry-of-life:13}
الملح: $9 \times 0.5 = \qty{4.5}{g}$؛ والغلوكوز: $1 \times 0.5 =
\qty{0.5}{g}$. أما الماء النقي فيخفف البلازما؛ فيدخل الماء عندئذ إلى
الكريات الحمراء، فتنتفخ وتنفجر.
\end{solution}

\begin{solution}{exo:g10:chemistry-of-life:14}
الغليكوجين: $500 \times 17 = \qty{8500}{kJ}$. والدهن: $10000 \times 38 =
\qty{3.8e5}{kJ}$. ولتخزين \qty{3.8e5}{kJ} غليكوجينًا: $380000/17
\approx \qty{22.4}{kg}$ من الغليكوجين، مع ثلاثة أمثال تلك الكتلة ماءً،
أي نحو \qty{90}{kg} إجمالًا --- وهو أكثر من كتلة الشخص نفسه. فالشحم
يخزن نحو تسعة أمثال الطاقة لكل كيلوغرام محمول، ولذلك تخزن الحيوانات
التي عليها أن تتحرك مدخراتها دهنًا.
\end{solution}

\begin{solution}{exo:g10:chemistry-of-life:15}
الذرات عادية بالفعل، لكن (1) النسب ليست كذلك: فالمادة الحية تركز
الكربون ألف ضعف والنيتروجين أكثر من ذلك بكثير، في مقابل قشرة من
الأكسجين والسيليسيوم؛ (2) والجزيئات ليست كذلك: فالكائنات الحية وحدها
تركب الكربون في سلاسل عملاقة --- بروتينات ونشاء وDNA --- لا نظير
لحجمها في المعادن؛ (3) وتلك الجزيئات تحمل معلومة في ترتيب وحداتها،
وما من بلورة تفعل ذلك. فالحياة خاصة كيميائيًا بتنظيمها، لا بذراتها.
\end{solution}

\begin{solution}{pb:g10:chemistry-of-life:1}
\textbf{1.} الماء $0.13 \times 45 = \qty{5.9}{mg}$؛ والكتلة الجافة
\qty{39.2}{mg}.

\textbf{2.} النشاء $0.70 \times 39.2 \approx \qty{27.4}{mg}$؛ \omterm{def:g10:chemistry-of-life:families}{والبروتين}
$0.12 \times 39.2 \approx \qty{4.7}{mg}$؛ والدهن $0.02 \times 39.2
\approx \qty{0.8}{mg}$.

\textbf{3.} اليود (أزرق أسود) للنشاء؛ والبيوريه (بنفسجي) \omterm{def:g10:chemistry-of-life:families}{للبروتين}؛
وبقعة شفافة على ورق أو السودان III للدهن.

\textbf{4.} عند 13\% تكون المدخرات صلبة والإنزيمات خاملة؛ وعلى الماء
أن يحل المدخرات ويطلق التفاعلات قبل أن يقدر الرشيم على النمو. فالإنبات
يبدأ بامتصاص الماء، ونسبة 88\% في البادرة علامة على أن كيمياءها تعمل.

\textbf{5.} الماء $0.13 \times 8 = \qty{1.04}{t}$؛ والمادة الجافة
\qty{6.96}{t}، أي نحو \qty{7}{t}.

\textbf{6.} الكتلة $500 + 320 + 10 - 130 = \qty{700}{g}$. والماء في
الرغيف: من الدقيق $0.13 \times 500 = \qty{65}{g}$، زائد \qty{320}{g}،
ناقص \qty{130}{g}: أي \qty{255}{g}، وهو $255/700 \approx 36\%$.

\textbf{7.} المادة الجافة في الدقيق $0.87 \times 500 = \qty{435}{g}$؛
والنشاء $0.70 \times 435 \approx \qty{305}{g}$؛ \omterm{def:g10:chemistry-of-life:families}{والبروتين} $0.12 \times 435 \approx
\qty{52}{g}$ (والدهن نحو \qty{9}{g}).

\textbf{8.} $305 \times 17 + 52 \times 17 + 9 \times 38 \approx 5185 +
884 + 342 \approx \qty{6400}{kJ}$.

\textbf{9.} $120/700 \approx 0.17$ من الرغيف: أي نحو \qty{1100}{kJ}،
وهو 11\% من \qty{10000}{kJ}.

\textbf{10.} الملح معدني: فهو مؤكسد تمامًا من قبل وليس فيه روابط
كربون--هيدروجين تحترق. أما في الجسم فشوارده تضبط ملوحة الدم ولا بد
منها لإشارات العصب والعضلة.

\textbf{11.} $0.60 \times 60 = \qty{36}{kg}$ من الماء.

\textbf{12.} \omterm{def:g10:chemistry-of-life:families}{البروتين} $0.17 \times 60 = \qty{10.2}{kg}$؛ والدهن $0.16
\times 60 = \qty{9.6}{kg}$؛ والمعادن \qty{3.6}{kg}؛ \omterm{def:g10:chemistry-of-life:families}{والكربوهيدرات}
\qty{0.6}{kg}.

\textbf{13.} $1.2/36 \approx 3.3\%$ من مائها. ويتوقف على ذلك الماء حجم
الدم وضبط الحرارة معًا؛ وفقدان 10\% خطر، فعلى الشرب أن يعيده في غضون
ساعات.

\textbf{14.} البلازما $0.55 \times 4.5 \approx \qty{2.5}{L}$؛ والماء $0.9
\times 2.5 \approx \qty{2.2}{kg}$، أي نحو 6\% من ماء الجسم. أما الباقي
فداخل الخلايا (نحو ثلثي ماء الجسم كله) وفي السائل بين الخلايا.

\textbf{15.} يتبادل الرضيع كل يوم نسبة من مائه أكبر بكثير مما يتبادله
البالغ (شربًا وبولًا وعبر الجلد)، فبضع ساعات من الفقد بلا وارد تسحب
منه نسبة خطرة من مجموعه.

\textbf{16.} \omterm{def:g10:chemistry-of-life:families}{الكربوهيدرات} $0.40 \times 0.6 = \qty{0.24}{kg}$؛ \omterm{def:g10:chemistry-of-life:families}{والدهون}
$0.77 \times 9.6 \approx \qty{7.4}{kg}$؛ \omterm{def:g10:chemistry-of-life:families}{والبروتينات} $0.53 \times 10.2
\approx \qty{5.4}{kg}$. والمجموع نحو \qty{13}{kg} من الكربون.

\textbf{17.} $13/60 \approx 22\%$، مقابل 18.5\%. والتوافق حسن بالنظر
إلى تقريب نسب العائلات وإلى تفاوت محتوى الدهن بين الأفراد؛ فالجسم
الأنحل يعطي رقمًا أدنى.

\textbf{18.} $0.16 \times 10.2 \approx \qty{1.6}{kg}$ من النيتروجين.
وقد جاء من نيتروجين الهواء، التقطته بكتيريا التربة، وامتصته النباتات،
وانتقل بالأكل على طول السلسلة الغذائية.

\textbf{19.} $0.89 \times 36 \approx \qty{32}{kg}$ من الأكسجين، أي
$32/60 \approx 53\%$ من كتلة الجسم. فالماء وحده يفسر معظم نسبة
65\%؛ والباقي هو الأكسجين داخل الجزيئات العضوية.

\textbf{20.} يحمل شخص كتلته \qty{60}{kg} نحو \qty{13}{kg} من الكربون.
وقد غادر الهواء ثاني أكسيد كربون بالتركيب الضوئي في النباتات، ودخل
الجسم بالتغذي والهضم على طول السلسلة الغذائية.
\end{solution}
